% Research Article: journal-style paper without a venue class. Pair with Minimal Research.
% When a venue is named, use its official class instead and keep only the structure.
\documentclass[11pt,a4paper]{article}

\usepackage[english]{babel}
\usepackage[T1]{fontenc}
\usepackage{microtype}
\usepackage[a4paper,margin=27mm]{geometry}
\usepackage{amsmath,mathtools,amsthm} % amsthm before newtxmath
\usepackage{newtxtext,newtxmath}
\usepackage{booktabs}
\usepackage{siunitx}
\usepackage{graphicx}
\usepackage[font=small,labelfont=bf]{caption}
\usepackage{subcaption}
\usepackage[hidelinks]{hyperref}
\usepackage[capitalise,noabbrev]{cleveref}

\newtheorem{proposition}{Proposition}
\widowpenalty=10000 \clubpenalty=10000

\title{\vspace{-2em}[CONTENT: title stating the finding]}
\author{[CONTENT: authors]\\[2pt]{\small [CONTENT: affiliations]}}
\date{}

\begin{document}
\maketitle

\begin{abstract}
\noindent % [CONTENT: problem, approach, main result with its number, implication; 150--250 words]
{}[CONTENT: abstract]
\end{abstract}

\noindent\textbf{Keywords:} [CONTENT: 3--6 keywords]

\section{Introduction}\label{sec:intro}
% Problem and why it matters; gap in prior work; contribution stated plainly; outline.
{}[CONTENT: introduction]

\section{Related work}\label{sec:related}
{}[CONTENT: positioned against verified sources only]

\section{Method}\label{sec:method}
Define every symbol at first use.
\begin{equation}\label{eq:model}
  y_i = \beta_0 + \beta_1 x_i + \varepsilon_i, \qquad \varepsilon_i \sim \mathcal{N}(0,\sigma^2).
\end{equation}

\section{Results}\label{sec:results}
\begin{table}[ht]
\centering
\caption{[CONTENT: what the reader should see in this table]}\label{tab:main}
\begin{tabular}{l S[table-format=2.2] S[table-format=1.2] S[table-format=3.0]}
\toprule
{Condition} & {Mean} & {SD} & {$n$} \\
\midrule
{}[CONTENT] & 0.00 & 0.00 & 0 \\
\bottomrule
\end{tabular}
\end{table}

\begin{figure}[ht]
\centering
\begin{subfigure}{.48\linewidth}\centering
  \fbox{\parbox[c][3cm][c]{.9\linewidth}{\centering\small[CONTENT: panel a, vector PDF]}}
  \caption{[CONTENT]}\label{fig:a}
\end{subfigure}\hfill
\begin{subfigure}{.48\linewidth}\centering
  \fbox{\parbox[c][3cm][c]{.9\linewidth}{\centering\small[CONTENT: panel b]}}
  \caption{[CONTENT]}\label{fig:b}
\end{subfigure}
\caption{[CONTENT: caption that states the takeaway]}\label{fig:main}
\end{figure}

As \cref{tab:main} and \cref{fig:main} show, [CONTENT].

\section{Discussion}\label{sec:discussion}
{}[CONTENT: interpretation, limitations, threats to validity]

\section{Conclusion}
{}[CONTENT]

\paragraph{Data and code availability.} [CONTENT]

\begin{thebibliography}{9}
% Only verified entries. Never write an entry from memory.
\bibitem{key} [CONTENT: verified reference]
\end{thebibliography}

\end{document}
